% Generated by `taskgen course examples` from crypto.rsa_toy (seed 3).
% Do not edit: change examples.yaml and regenerate.
\begin{example}[ключи RSA, шифрование и расшифрование по КТО]\label{ex:rsa-toy}
Ключи RSA строятся по простым $p = 23$, $q = 17$ и открытой экспоненте $e = 11$ (если $\gcd(e, \varphi(N)) \ne 1$, вместо $e$ берётся наименьшее большее его число, взаимно простое с $\varphi(N)$). Найти $N$, $\varphi(N)$ и закрытый ключ $d$ (расширенным алгоритмом Евклида). Зашифровать сообщение $m = 98$ бинарным алгоритмом возведения в степень и расшифровать полученный шифротекст с помощью китайской теоремы об остатках.

\par\noindent\textit{Решение.} $N = pq = 23 \cdot 17 = 391$, $\varphi(N) = (p - 1)(q - 1) = 22 \cdot 16 = 352 = 2^{5} \cdot 11$. $\gcd(11, 352) = 11 \ne 1$: такой показатель недопустим (обратного к $e$ по модулю $\varphi(N)$ нет, и отображение $x \mapsto x^e$ не биективно). Перебираем следующие числа: $\gcd(12, 352) = 4$, $\gcd(13, 352) = 1$, поэтому $e = 13$.\par Закрытый ключ $d \equiv e^{-1} \pmod{\varphi(N)}$.  Обратный к $13$ по модулю $352$ (расширенный алгоритм Евклида для $(352, 13)$):
\[\begin{array}{r|r|r|r}
r_i & q_i & u_i & v_i \\ \hline
352 &  & 1 & 0 \\
13 &  & 0 & 1 \\
1 & 27 & 1 & -27 \\
0 & 13 & -13 & 352 \\
\end{array}\qquad r_i = u_i \cdot 352 + v_i \cdot 13\]

$13^{-1} \equiv -27 \equiv 325\ (\textrm{mod}\ 352)$.
 Проверка: $ed = 13 \cdot 325 = 4225 = 12 \cdot 352 + 1$.\par Шифрование: $c \equiv m^e \pmod{N}$, $e = 13 = 1101_2$, все вычисления по модулю $391$:\[\begin{array}{c|r|r|r}\text{бит} & r^2 & r & e \\ \hline 1 & 1 & 98 & 1 \\ 1 & 220 & 55 & 3 \\ 0 & 288 & 288 & 6 \\ 1 & 52 & 13 & 13 \\ \end{array}\qquad c = 13.\]Расшифрование по КТО: $d_p = d \bmod (p - 1) = 17$, $d_q = d \bmod (q - 1) = 5$, $c_p = c \bmod p = 13$, $c_q = c \bmod q = 13$; $m_p = c_p^{d_p} \bmod p$, $m_q = c_q^{d_q} \bmod q$.\[\begin{array}{c|r|r|r}\text{бит} & r^2 & r & e \\ \hline 1 & 1 & 13 & 1 \\ 0 & 8 & 8 & 2 \\ 0 & 18 & 18 & 4 \\ 0 & 2 & 2 & 8 \\ 1 & 4 & 6 & 17 \\ \end{array}\ \ m_p = 6\qquad \begin{array}{c|r|r|r}\text{бит} & r^2 & r & e \\ \hline 1 & 1 & 13 & 1 \\ 0 & 16 & 16 & 2 \\ 1 & 1 & 13 & 5 \\ \end{array}\ \ m_q = 13\]\[\left\{\begin{array}{l}
x \equiv 6\ (\textrm{mod}\ 23) \\
x \equiv 13\ (\textrm{mod}\ 17) \\
\end{array}\right.\qquad M = 391.\]
$M_{1} = M/23 = 17 \equiv 17$, $y_{1} = M_{1}^{-1} \equiv 19\ (\textrm{mod}\ 23)$; $M_{2} = M/17 = 23 \equiv 6$, $y_{2} = M_{2}^{-1} \equiv 3\ (\textrm{mod}\ 17)$.\[\begin{aligned}x &\equiv \textstyle\sum_i r_i M_i y_i \\ &= 6\cdot17\cdot19 + 13\cdot23\cdot3 \\ &\equiv 98\ (\textrm{mod}\ 391).\end{aligned}\]
\[N = 391,\ \varphi(N) = 352,\ e = 13,\ d = 325,\ c = 13,\ m = 98.\]
\end{example}
